% ═══════════════════════════════════════════════════════════════
% 民事起诉状 · 共享前导码
% 所有民事起诉状模板共用本文件，通过 % ═══════════════════════════════════════════════════════════════
% 民事起诉状 · 共享前导码
% 所有民事起诉状模板共用本文件，通过 % ═══════════════════════════════════════════════════════════════
% 民事起诉状 · 共享前导码
% 所有民事起诉状模板共用本文件，通过 % ═══════════════════════════════════════════════════════════════
% 民事起诉状 · 共享前导码
% 所有民事起诉状模板共用本文件，通过 \input{../common/preamble} 引入
% ═══════════════════════════════════════════════════════════════

\documentclass[a4paper,12pt,twoside]{ctexart}
\usepackage[dvipsnames]{xcolor}
\usepackage{geometry}
\usepackage{tcolorbox}
\usepackage{tikz}
\usepackage{amssymb}
\usepackage{graphicx}
\usepackage{tabularx}
\usepackage{makecell}
\usepackage{multirow}

% 定义公函标准炭黑（K=92%，RGB约0.08）：权威稳重、高质量印刷感
\definecolor{inkblack}{rgb}{0.08,0.08,0.08}

\geometry{a4paper, left=3.18cm, right=3.18cm, top=2.54cm, bottom=2.54cm}

% 配置页脚，实现奇数页右下角，偶数页左下角的页码
\usepackage{fancyhdr}
\pagestyle{fancy}
\fancyhf{}
\renewcommand{\headrulewidth}{0pt}

% 自定义页脚位置微调
\fancyfoot[RO]{\makebox[0pt][l]{\hspace*{28.2pt}\raisebox{2.81pt}{\rmfamily\fontsize{10.5bp}{10.5bp}\selectfont\thepage}}}
\fancyfoot[LE]{\makebox[0pt][r]{\hspace*{33.45pt}\raisebox{2.81pt}{\rmfamily\fontsize{10.5bp}{10.5bp}\selectfont\thepage}}}

% 字体声明
% \fontpath 可由各模板在 \input{preamble} 之前通过
% \newcommand{\fontpath}{...} 覆盖设置
% 默认路径：从 templates/XXX/ 编译时指向 templates/common/fonts/
\providecommand{\fontpath}{../common/fonts/}
\newCJKfontfamily\fzdbs[Path=\fontpath]{FZDBS.ttf}
\newCJKfontfamily\fzssk[Path=\fontpath]{FZSSK.ttf}

% 全局文档颜色设定：公函炭黑
\AtBeginDocument{\color{inkblack}}

% 全局禁止段间距拉伸，确保 tcolorbox 框之间无缝隙
\setlength{\parskip}{0pt}

% ─── 复合标题样式 ───
\newcommand{\lawtitle}[2]{
    \vspace*{0.507cm} 
    \begin{center}
        {\fzdbs\fontsize{22bp}{22bp}\selectfont #1}\par
        \vspace{10.37bp}
        {\fzdbs\fontsize{18bp}{18bp}\selectfont #2}\par
    \end{center}
    \vspace{18.4bp}
}

% ─── 表框（说明框）环境 ───
\newtcolorbox{notesection}[1][]{
    width=467.72bp,
    sharp corners,
    colback=white,
    colframe=inkblack,
    boxrule=0.5pt,
    fontupper=\color{inkblack}\fontsize{10.5bp}{10.5bp}\selectfont,
    enlarge left by=-19.28bp,
    before={\par\nointerlineskip\vspace{-0.5pt}\noindent},
    after={\par},
    boxsep=0pt,
    left=4bp,
    top=4.69bp,
    bottom=4.69bp,
    right=0pt,
    #1
}

% ─── 法律段落换行命令 ───
\newcommand{\lawpar}{\par\vspace*{-4.86bp}}

% ─── 区块副标题样式 ───
\newcommand{\subtitlesection}[1]{%
    \begin{tcolorbox}[
        width=467.72bp,
        sharp corners,
        colback=white,
        colframe=inkblack,
        boxrule=0.5pt,
        boxsep=0pt,
        left=0pt, right=0pt,
        top=4bp, bottom=4bp,
        halign=center,
        fontupper=\fzdbs\fontsize{15bp}{15bp}\selectfont\color{inkblack},
        enlarge left by=-19.28bp,
        before={\par\nointerlineskip\vspace{-0.5pt}\noindent},
        after={\par}
    ]
    #1
    \end{tcolorbox}%
}

% ─── 区块副标题样式（带说明文字） ───
\newcommand{\subtitlenotesection}[2]{%
    \begin{tcolorbox}[
        width=467.72bp, sharp corners, colback=white, colframe=inkblack, boxrule=0.5pt,
        boxsep=0pt, left=0pt, right=0pt, top=4bp, bottom=4bp, halign=center,
        enlarge left by=-19.28bp,
        before={\par\nointerlineskip\vspace{-0.5pt}\noindent}, after={\par}
    ]
    {\fzdbs\fontsize{15bp}{15bp}\selectfont\color{inkblack} #1}\par\nointerlineskip
    \vspace{2bp}
    {\heiti\fontsize{10.5bp}{14bp}\selectfont\color{inkblack} #2}
    \end{tcolorbox}%
}

% ─── 左右分栏基础命令 ───
\newcommand{\pairsectionbase}[3]{%
    \begin{tcolorbox}[
        width=467.72bp, sharp corners, colback=white, colframe=inkblack, boxrule=0.5pt,
        enlarge left by=-19.28bp,
        before={\par\nointerlineskip\vspace{-0.5pt}\noindent}, after={\par},
        boxsep=0pt, left=0pt, right=0pt, top=0pt, bottom=0pt
    ]
    \begin{minipage}[c]{112.9bp}%
        \vspace{2bp}%
        \songti\fontsize{10.5bp}{12bp}\selectfont
        #1#2%
        \vspace{2bp}%
    \end{minipage}%
    \hspace{0pt}%
    \vrule width 0.5pt%
    \hspace{4bp}%
    \begin{minipage}[c]{346.32bp}%
        \vspace{2bp}%
        \songti\mdseries\fontsize{10.5bp}{17bp}\selectfont%
        \baselineskip=17bp%
        \setlength{\parindent}{0pt}%
        \color{inkblack}%
        #3%
        \vspace{2bp}%
    \end{minipage}%
    \end{tcolorbox}%
}

% 标签居中版
\newcommand{\pairsection}[2]{%
    \pairsectionbase{\centering}{#1}{#2}%
}

% 标签左对齐版（用于长标签）
\newcommand{\pairsectionleft}[2]{%
    \pairsectionbase{\raggedright\fontsize{10.5bp}{10.5bp}\selectfont}{#1}{#2}%
}


 引入
% ═══════════════════════════════════════════════════════════════

\documentclass[a4paper,12pt,twoside]{ctexart}
\usepackage[dvipsnames]{xcolor}
\usepackage{geometry}
\usepackage{tcolorbox}
\usepackage{tikz}
\usepackage{amssymb}
\usepackage{graphicx}
\usepackage{tabularx}
\usepackage{makecell}
\usepackage{multirow}

% 定义公函标准炭黑（K=92%，RGB约0.08）：权威稳重、高质量印刷感
\definecolor{inkblack}{rgb}{0.08,0.08,0.08}

\geometry{a4paper, left=3.18cm, right=3.18cm, top=2.54cm, bottom=2.54cm}

% 配置页脚，实现奇数页右下角，偶数页左下角的页码
\usepackage{fancyhdr}
\pagestyle{fancy}
\fancyhf{}
\renewcommand{\headrulewidth}{0pt}

% 自定义页脚位置微调
\fancyfoot[RO]{\makebox[0pt][l]{\hspace*{28.2pt}\raisebox{2.81pt}{\rmfamily\fontsize{10.5bp}{10.5bp}\selectfont\thepage}}}
\fancyfoot[LE]{\makebox[0pt][r]{\hspace*{33.45pt}\raisebox{2.81pt}{\rmfamily\fontsize{10.5bp}{10.5bp}\selectfont\thepage}}}

% 字体声明
% \fontpath 可由各模板在 % ═══════════════════════════════════════════════════════════════
% 民事起诉状 · 共享前导码
% 所有民事起诉状模板共用本文件，通过 \input{../common/preamble} 引入
% ═══════════════════════════════════════════════════════════════

\documentclass[a4paper,12pt,twoside]{ctexart}
\usepackage[dvipsnames]{xcolor}
\usepackage{geometry}
\usepackage{tcolorbox}
\usepackage{tikz}
\usepackage{amssymb}
\usepackage{graphicx}
\usepackage{tabularx}
\usepackage{makecell}
\usepackage{multirow}

% 定义公函标准炭黑（K=92%，RGB约0.08）：权威稳重、高质量印刷感
\definecolor{inkblack}{rgb}{0.08,0.08,0.08}

\geometry{a4paper, left=3.18cm, right=3.18cm, top=2.54cm, bottom=2.54cm}

% 配置页脚，实现奇数页右下角，偶数页左下角的页码
\usepackage{fancyhdr}
\pagestyle{fancy}
\fancyhf{}
\renewcommand{\headrulewidth}{0pt}

% 自定义页脚位置微调
\fancyfoot[RO]{\makebox[0pt][l]{\hspace*{28.2pt}\raisebox{2.81pt}{\rmfamily\fontsize{10.5bp}{10.5bp}\selectfont\thepage}}}
\fancyfoot[LE]{\makebox[0pt][r]{\hspace*{33.45pt}\raisebox{2.81pt}{\rmfamily\fontsize{10.5bp}{10.5bp}\selectfont\thepage}}}

% 字体声明
% \fontpath 可由各模板在 \input{preamble} 之前通过
% \newcommand{\fontpath}{...} 覆盖设置
% 默认路径：从 templates/XXX/ 编译时指向 templates/common/fonts/
\providecommand{\fontpath}{../common/fonts/}
\newCJKfontfamily\fzdbs[Path=\fontpath]{FZDBS.ttf}
\newCJKfontfamily\fzssk[Path=\fontpath]{FZSSK.ttf}

% 全局文档颜色设定：公函炭黑
\AtBeginDocument{\color{inkblack}}

% 全局禁止段间距拉伸，确保 tcolorbox 框之间无缝隙
\setlength{\parskip}{0pt}

% ─── 复合标题样式 ───
\newcommand{\lawtitle}[2]{
    \vspace*{0.507cm} 
    \begin{center}
        {\fzdbs\fontsize{22bp}{22bp}\selectfont #1}\par
        \vspace{10.37bp}
        {\fzdbs\fontsize{18bp}{18bp}\selectfont #2}\par
    \end{center}
    \vspace{18.4bp}
}

% ─── 表框（说明框）环境 ───
\newtcolorbox{notesection}[1][]{
    width=467.72bp,
    sharp corners,
    colback=white,
    colframe=inkblack,
    boxrule=0.5pt,
    fontupper=\color{inkblack}\fontsize{10.5bp}{10.5bp}\selectfont,
    enlarge left by=-19.28bp,
    before={\par\nointerlineskip\vspace{-0.5pt}\noindent},
    after={\par},
    boxsep=0pt,
    left=4bp,
    top=4.69bp,
    bottom=4.69bp,
    right=0pt,
    #1
}

% ─── 法律段落换行命令 ───
\newcommand{\lawpar}{\par\vspace*{-4.86bp}}

% ─── 区块副标题样式 ───
\newcommand{\subtitlesection}[1]{%
    \begin{tcolorbox}[
        width=467.72bp,
        sharp corners,
        colback=white,
        colframe=inkblack,
        boxrule=0.5pt,
        boxsep=0pt,
        left=0pt, right=0pt,
        top=4bp, bottom=4bp,
        halign=center,
        fontupper=\fzdbs\fontsize{15bp}{15bp}\selectfont\color{inkblack},
        enlarge left by=-19.28bp,
        before={\par\nointerlineskip\vspace{-0.5pt}\noindent},
        after={\par}
    ]
    #1
    \end{tcolorbox}%
}

% ─── 区块副标题样式（带说明文字） ───
\newcommand{\subtitlenotesection}[2]{%
    \begin{tcolorbox}[
        width=467.72bp, sharp corners, colback=white, colframe=inkblack, boxrule=0.5pt,
        boxsep=0pt, left=0pt, right=0pt, top=4bp, bottom=4bp, halign=center,
        enlarge left by=-19.28bp,
        before={\par\nointerlineskip\vspace{-0.5pt}\noindent}, after={\par}
    ]
    {\fzdbs\fontsize{15bp}{15bp}\selectfont\color{inkblack} #1}\par\nointerlineskip
    \vspace{2bp}
    {\heiti\fontsize{10.5bp}{14bp}\selectfont\color{inkblack} #2}
    \end{tcolorbox}%
}

% ─── 左右分栏基础命令 ───
\newcommand{\pairsectionbase}[3]{%
    \begin{tcolorbox}[
        width=467.72bp, sharp corners, colback=white, colframe=inkblack, boxrule=0.5pt,
        enlarge left by=-19.28bp,
        before={\par\nointerlineskip\vspace{-0.5pt}\noindent}, after={\par},
        boxsep=0pt, left=0pt, right=0pt, top=0pt, bottom=0pt
    ]
    \begin{minipage}[c]{112.9bp}%
        \vspace{2bp}%
        \songti\fontsize{10.5bp}{12bp}\selectfont
        #1#2%
        \vspace{2bp}%
    \end{minipage}%
    \hspace{0pt}%
    \vrule width 0.5pt%
    \hspace{4bp}%
    \begin{minipage}[c]{346.32bp}%
        \vspace{2bp}%
        \songti\mdseries\fontsize{10.5bp}{17bp}\selectfont%
        \baselineskip=17bp%
        \setlength{\parindent}{0pt}%
        \color{inkblack}%
        #3%
        \vspace{2bp}%
    \end{minipage}%
    \end{tcolorbox}%
}

% 标签居中版
\newcommand{\pairsection}[2]{%
    \pairsectionbase{\centering}{#1}{#2}%
}

% 标签左对齐版（用于长标签）
\newcommand{\pairsectionleft}[2]{%
    \pairsectionbase{\raggedright\fontsize{10.5bp}{10.5bp}\selectfont}{#1}{#2}%
}


 之前通过
% \newcommand{\fontpath}{...} 覆盖设置
% 默认路径：从 templates/XXX/ 编译时指向 templates/common/fonts/
\providecommand{\fontpath}{../common/fonts/}
\newCJKfontfamily\fzdbs[Path=\fontpath]{FZDBS.ttf}
\newCJKfontfamily\fzssk[Path=\fontpath]{FZSSK.ttf}

% 全局文档颜色设定：公函炭黑
\AtBeginDocument{\color{inkblack}}

% 全局禁止段间距拉伸，确保 tcolorbox 框之间无缝隙
\setlength{\parskip}{0pt}

% ─── 复合标题样式 ───
\newcommand{\lawtitle}[2]{
    \vspace*{0.507cm} 
    \begin{center}
        {\fzdbs\fontsize{22bp}{22bp}\selectfont #1}\par
        \vspace{10.37bp}
        {\fzdbs\fontsize{18bp}{18bp}\selectfont #2}\par
    \end{center}
    \vspace{18.4bp}
}

% ─── 表框（说明框）环境 ───
\newtcolorbox{notesection}[1][]{
    width=467.72bp,
    sharp corners,
    colback=white,
    colframe=inkblack,
    boxrule=0.5pt,
    fontupper=\color{inkblack}\fontsize{10.5bp}{10.5bp}\selectfont,
    enlarge left by=-19.28bp,
    before={\par\nointerlineskip\vspace{-0.5pt}\noindent},
    after={\par},
    boxsep=0pt,
    left=4bp,
    top=4.69bp,
    bottom=4.69bp,
    right=0pt,
    #1
}

% ─── 法律段落换行命令 ───
\newcommand{\lawpar}{\par\vspace*{-4.86bp}}

% ─── 区块副标题样式 ───
\newcommand{\subtitlesection}[1]{%
    \begin{tcolorbox}[
        width=467.72bp,
        sharp corners,
        colback=white,
        colframe=inkblack,
        boxrule=0.5pt,
        boxsep=0pt,
        left=0pt, right=0pt,
        top=4bp, bottom=4bp,
        halign=center,
        fontupper=\fzdbs\fontsize{15bp}{15bp}\selectfont\color{inkblack},
        enlarge left by=-19.28bp,
        before={\par\nointerlineskip\vspace{-0.5pt}\noindent},
        after={\par}
    ]
    #1
    \end{tcolorbox}%
}

% ─── 区块副标题样式（带说明文字） ───
\newcommand{\subtitlenotesection}[2]{%
    \begin{tcolorbox}[
        width=467.72bp, sharp corners, colback=white, colframe=inkblack, boxrule=0.5pt,
        boxsep=0pt, left=0pt, right=0pt, top=4bp, bottom=4bp, halign=center,
        enlarge left by=-19.28bp,
        before={\par\nointerlineskip\vspace{-0.5pt}\noindent}, after={\par}
    ]
    {\fzdbs\fontsize{15bp}{15bp}\selectfont\color{inkblack} #1}\par\nointerlineskip
    \vspace{2bp}
    {\heiti\fontsize{10.5bp}{14bp}\selectfont\color{inkblack} #2}
    \end{tcolorbox}%
}

% ─── 左右分栏基础命令 ───
\newcommand{\pairsectionbase}[3]{%
    \begin{tcolorbox}[
        width=467.72bp, sharp corners, colback=white, colframe=inkblack, boxrule=0.5pt,
        enlarge left by=-19.28bp,
        before={\par\nointerlineskip\vspace{-0.5pt}\noindent}, after={\par},
        boxsep=0pt, left=0pt, right=0pt, top=0pt, bottom=0pt
    ]
    \begin{minipage}[c]{112.9bp}%
        \vspace{2bp}%
        \songti\fontsize{10.5bp}{12bp}\selectfont
        #1#2%
        \vspace{2bp}%
    \end{minipage}%
    \hspace{0pt}%
    \vrule width 0.5pt%
    \hspace{4bp}%
    \begin{minipage}[c]{346.32bp}%
        \vspace{2bp}%
        \songti\mdseries\fontsize{10.5bp}{17bp}\selectfont%
        \baselineskip=17bp%
        \setlength{\parindent}{0pt}%
        \color{inkblack}%
        #3%
        \vspace{2bp}%
    \end{minipage}%
    \end{tcolorbox}%
}

% 标签居中版
\newcommand{\pairsection}[2]{%
    \pairsectionbase{\centering}{#1}{#2}%
}

% 标签左对齐版（用于长标签）
\newcommand{\pairsectionleft}[2]{%
    \pairsectionbase{\raggedright\fontsize{10.5bp}{10.5bp}\selectfont}{#1}{#2}%
}


 引入
% ═══════════════════════════════════════════════════════════════

\documentclass[a4paper,12pt,twoside]{ctexart}
\usepackage[dvipsnames]{xcolor}
\usepackage{geometry}
\usepackage{tcolorbox}
\usepackage{tikz}
\usepackage{amssymb}
\usepackage{graphicx}
\usepackage{tabularx}
\usepackage{makecell}
\usepackage{multirow}

% 定义公函标准炭黑（K=92%，RGB约0.08）：权威稳重、高质量印刷感
\definecolor{inkblack}{rgb}{0.08,0.08,0.08}

\geometry{a4paper, left=3.18cm, right=3.18cm, top=2.54cm, bottom=2.54cm}

% 配置页脚，实现奇数页右下角，偶数页左下角的页码
\usepackage{fancyhdr}
\pagestyle{fancy}
\fancyhf{}
\renewcommand{\headrulewidth}{0pt}

% 自定义页脚位置微调
\fancyfoot[RO]{\makebox[0pt][l]{\hspace*{28.2pt}\raisebox{2.81pt}{\rmfamily\fontsize{10.5bp}{10.5bp}\selectfont\thepage}}}
\fancyfoot[LE]{\makebox[0pt][r]{\hspace*{33.45pt}\raisebox{2.81pt}{\rmfamily\fontsize{10.5bp}{10.5bp}\selectfont\thepage}}}

% 字体声明
% \fontpath 可由各模板在 % ═══════════════════════════════════════════════════════════════
% 民事起诉状 · 共享前导码
% 所有民事起诉状模板共用本文件，通过 % ═══════════════════════════════════════════════════════════════
% 民事起诉状 · 共享前导码
% 所有民事起诉状模板共用本文件，通过 \input{../common/preamble} 引入
% ═══════════════════════════════════════════════════════════════

\documentclass[a4paper,12pt,twoside]{ctexart}
\usepackage[dvipsnames]{xcolor}
\usepackage{geometry}
\usepackage{tcolorbox}
\usepackage{tikz}
\usepackage{amssymb}
\usepackage{graphicx}
\usepackage{tabularx}
\usepackage{makecell}
\usepackage{multirow}

% 定义公函标准炭黑（K=92%，RGB约0.08）：权威稳重、高质量印刷感
\definecolor{inkblack}{rgb}{0.08,0.08,0.08}

\geometry{a4paper, left=3.18cm, right=3.18cm, top=2.54cm, bottom=2.54cm}

% 配置页脚，实现奇数页右下角，偶数页左下角的页码
\usepackage{fancyhdr}
\pagestyle{fancy}
\fancyhf{}
\renewcommand{\headrulewidth}{0pt}

% 自定义页脚位置微调
\fancyfoot[RO]{\makebox[0pt][l]{\hspace*{28.2pt}\raisebox{2.81pt}{\rmfamily\fontsize{10.5bp}{10.5bp}\selectfont\thepage}}}
\fancyfoot[LE]{\makebox[0pt][r]{\hspace*{33.45pt}\raisebox{2.81pt}{\rmfamily\fontsize{10.5bp}{10.5bp}\selectfont\thepage}}}

% 字体声明
% \fontpath 可由各模板在 \input{preamble} 之前通过
% \newcommand{\fontpath}{...} 覆盖设置
% 默认路径：从 templates/XXX/ 编译时指向 templates/common/fonts/
\providecommand{\fontpath}{../common/fonts/}
\newCJKfontfamily\fzdbs[Path=\fontpath]{FZDBS.ttf}
\newCJKfontfamily\fzssk[Path=\fontpath]{FZSSK.ttf}

% 全局文档颜色设定：公函炭黑
\AtBeginDocument{\color{inkblack}}

% 全局禁止段间距拉伸，确保 tcolorbox 框之间无缝隙
\setlength{\parskip}{0pt}

% ─── 复合标题样式 ───
\newcommand{\lawtitle}[2]{
    \vspace*{0.507cm} 
    \begin{center}
        {\fzdbs\fontsize{22bp}{22bp}\selectfont #1}\par
        \vspace{10.37bp}
        {\fzdbs\fontsize{18bp}{18bp}\selectfont #2}\par
    \end{center}
    \vspace{18.4bp}
}

% ─── 表框（说明框）环境 ───
\newtcolorbox{notesection}[1][]{
    width=467.72bp,
    sharp corners,
    colback=white,
    colframe=inkblack,
    boxrule=0.5pt,
    fontupper=\color{inkblack}\fontsize{10.5bp}{10.5bp}\selectfont,
    enlarge left by=-19.28bp,
    before={\par\nointerlineskip\vspace{-0.5pt}\noindent},
    after={\par},
    boxsep=0pt,
    left=4bp,
    top=4.69bp,
    bottom=4.69bp,
    right=0pt,
    #1
}

% ─── 法律段落换行命令 ───
\newcommand{\lawpar}{\par\vspace*{-4.86bp}}

% ─── 区块副标题样式 ───
\newcommand{\subtitlesection}[1]{%
    \begin{tcolorbox}[
        width=467.72bp,
        sharp corners,
        colback=white,
        colframe=inkblack,
        boxrule=0.5pt,
        boxsep=0pt,
        left=0pt, right=0pt,
        top=4bp, bottom=4bp,
        halign=center,
        fontupper=\fzdbs\fontsize{15bp}{15bp}\selectfont\color{inkblack},
        enlarge left by=-19.28bp,
        before={\par\nointerlineskip\vspace{-0.5pt}\noindent},
        after={\par}
    ]
    #1
    \end{tcolorbox}%
}

% ─── 区块副标题样式（带说明文字） ───
\newcommand{\subtitlenotesection}[2]{%
    \begin{tcolorbox}[
        width=467.72bp, sharp corners, colback=white, colframe=inkblack, boxrule=0.5pt,
        boxsep=0pt, left=0pt, right=0pt, top=4bp, bottom=4bp, halign=center,
        enlarge left by=-19.28bp,
        before={\par\nointerlineskip\vspace{-0.5pt}\noindent}, after={\par}
    ]
    {\fzdbs\fontsize{15bp}{15bp}\selectfont\color{inkblack} #1}\par\nointerlineskip
    \vspace{2bp}
    {\heiti\fontsize{10.5bp}{14bp}\selectfont\color{inkblack} #2}
    \end{tcolorbox}%
}

% ─── 左右分栏基础命令 ───
\newcommand{\pairsectionbase}[3]{%
    \begin{tcolorbox}[
        width=467.72bp, sharp corners, colback=white, colframe=inkblack, boxrule=0.5pt,
        enlarge left by=-19.28bp,
        before={\par\nointerlineskip\vspace{-0.5pt}\noindent}, after={\par},
        boxsep=0pt, left=0pt, right=0pt, top=0pt, bottom=0pt
    ]
    \begin{minipage}[c]{112.9bp}%
        \vspace{2bp}%
        \songti\fontsize{10.5bp}{12bp}\selectfont
        #1#2%
        \vspace{2bp}%
    \end{minipage}%
    \hspace{0pt}%
    \vrule width 0.5pt%
    \hspace{4bp}%
    \begin{minipage}[c]{346.32bp}%
        \vspace{2bp}%
        \songti\mdseries\fontsize{10.5bp}{17bp}\selectfont%
        \baselineskip=17bp%
        \setlength{\parindent}{0pt}%
        \color{inkblack}%
        #3%
        \vspace{2bp}%
    \end{minipage}%
    \end{tcolorbox}%
}

% 标签居中版
\newcommand{\pairsection}[2]{%
    \pairsectionbase{\centering}{#1}{#2}%
}

% 标签左对齐版（用于长标签）
\newcommand{\pairsectionleft}[2]{%
    \pairsectionbase{\raggedright\fontsize{10.5bp}{10.5bp}\selectfont}{#1}{#2}%
}


 引入
% ═══════════════════════════════════════════════════════════════

\documentclass[a4paper,12pt,twoside]{ctexart}
\usepackage[dvipsnames]{xcolor}
\usepackage{geometry}
\usepackage{tcolorbox}
\usepackage{tikz}
\usepackage{amssymb}
\usepackage{graphicx}
\usepackage{tabularx}
\usepackage{makecell}
\usepackage{multirow}

% 定义公函标准炭黑（K=92%，RGB约0.08）：权威稳重、高质量印刷感
\definecolor{inkblack}{rgb}{0.08,0.08,0.08}

\geometry{a4paper, left=3.18cm, right=3.18cm, top=2.54cm, bottom=2.54cm}

% 配置页脚，实现奇数页右下角，偶数页左下角的页码
\usepackage{fancyhdr}
\pagestyle{fancy}
\fancyhf{}
\renewcommand{\headrulewidth}{0pt}

% 自定义页脚位置微调
\fancyfoot[RO]{\makebox[0pt][l]{\hspace*{28.2pt}\raisebox{2.81pt}{\rmfamily\fontsize{10.5bp}{10.5bp}\selectfont\thepage}}}
\fancyfoot[LE]{\makebox[0pt][r]{\hspace*{33.45pt}\raisebox{2.81pt}{\rmfamily\fontsize{10.5bp}{10.5bp}\selectfont\thepage}}}

% 字体声明
% \fontpath 可由各模板在 % ═══════════════════════════════════════════════════════════════
% 民事起诉状 · 共享前导码
% 所有民事起诉状模板共用本文件，通过 \input{../common/preamble} 引入
% ═══════════════════════════════════════════════════════════════

\documentclass[a4paper,12pt,twoside]{ctexart}
\usepackage[dvipsnames]{xcolor}
\usepackage{geometry}
\usepackage{tcolorbox}
\usepackage{tikz}
\usepackage{amssymb}
\usepackage{graphicx}
\usepackage{tabularx}
\usepackage{makecell}
\usepackage{multirow}

% 定义公函标准炭黑（K=92%，RGB约0.08）：权威稳重、高质量印刷感
\definecolor{inkblack}{rgb}{0.08,0.08,0.08}

\geometry{a4paper, left=3.18cm, right=3.18cm, top=2.54cm, bottom=2.54cm}

% 配置页脚，实现奇数页右下角，偶数页左下角的页码
\usepackage{fancyhdr}
\pagestyle{fancy}
\fancyhf{}
\renewcommand{\headrulewidth}{0pt}

% 自定义页脚位置微调
\fancyfoot[RO]{\makebox[0pt][l]{\hspace*{28.2pt}\raisebox{2.81pt}{\rmfamily\fontsize{10.5bp}{10.5bp}\selectfont\thepage}}}
\fancyfoot[LE]{\makebox[0pt][r]{\hspace*{33.45pt}\raisebox{2.81pt}{\rmfamily\fontsize{10.5bp}{10.5bp}\selectfont\thepage}}}

% 字体声明
% \fontpath 可由各模板在 \input{preamble} 之前通过
% \newcommand{\fontpath}{...} 覆盖设置
% 默认路径：从 templates/XXX/ 编译时指向 templates/common/fonts/
\providecommand{\fontpath}{../common/fonts/}
\newCJKfontfamily\fzdbs[Path=\fontpath]{FZDBS.ttf}
\newCJKfontfamily\fzssk[Path=\fontpath]{FZSSK.ttf}

% 全局文档颜色设定：公函炭黑
\AtBeginDocument{\color{inkblack}}

% 全局禁止段间距拉伸，确保 tcolorbox 框之间无缝隙
\setlength{\parskip}{0pt}

% ─── 复合标题样式 ───
\newcommand{\lawtitle}[2]{
    \vspace*{0.507cm} 
    \begin{center}
        {\fzdbs\fontsize{22bp}{22bp}\selectfont #1}\par
        \vspace{10.37bp}
        {\fzdbs\fontsize{18bp}{18bp}\selectfont #2}\par
    \end{center}
    \vspace{18.4bp}
}

% ─── 表框（说明框）环境 ───
\newtcolorbox{notesection}[1][]{
    width=467.72bp,
    sharp corners,
    colback=white,
    colframe=inkblack,
    boxrule=0.5pt,
    fontupper=\color{inkblack}\fontsize{10.5bp}{10.5bp}\selectfont,
    enlarge left by=-19.28bp,
    before={\par\nointerlineskip\vspace{-0.5pt}\noindent},
    after={\par},
    boxsep=0pt,
    left=4bp,
    top=4.69bp,
    bottom=4.69bp,
    right=0pt,
    #1
}

% ─── 法律段落换行命令 ───
\newcommand{\lawpar}{\par\vspace*{-4.86bp}}

% ─── 区块副标题样式 ───
\newcommand{\subtitlesection}[1]{%
    \begin{tcolorbox}[
        width=467.72bp,
        sharp corners,
        colback=white,
        colframe=inkblack,
        boxrule=0.5pt,
        boxsep=0pt,
        left=0pt, right=0pt,
        top=4bp, bottom=4bp,
        halign=center,
        fontupper=\fzdbs\fontsize{15bp}{15bp}\selectfont\color{inkblack},
        enlarge left by=-19.28bp,
        before={\par\nointerlineskip\vspace{-0.5pt}\noindent},
        after={\par}
    ]
    #1
    \end{tcolorbox}%
}

% ─── 区块副标题样式（带说明文字） ───
\newcommand{\subtitlenotesection}[2]{%
    \begin{tcolorbox}[
        width=467.72bp, sharp corners, colback=white, colframe=inkblack, boxrule=0.5pt,
        boxsep=0pt, left=0pt, right=0pt, top=4bp, bottom=4bp, halign=center,
        enlarge left by=-19.28bp,
        before={\par\nointerlineskip\vspace{-0.5pt}\noindent}, after={\par}
    ]
    {\fzdbs\fontsize{15bp}{15bp}\selectfont\color{inkblack} #1}\par\nointerlineskip
    \vspace{2bp}
    {\heiti\fontsize{10.5bp}{14bp}\selectfont\color{inkblack} #2}
    \end{tcolorbox}%
}

% ─── 左右分栏基础命令 ───
\newcommand{\pairsectionbase}[3]{%
    \begin{tcolorbox}[
        width=467.72bp, sharp corners, colback=white, colframe=inkblack, boxrule=0.5pt,
        enlarge left by=-19.28bp,
        before={\par\nointerlineskip\vspace{-0.5pt}\noindent}, after={\par},
        boxsep=0pt, left=0pt, right=0pt, top=0pt, bottom=0pt
    ]
    \begin{minipage}[c]{112.9bp}%
        \vspace{2bp}%
        \songti\fontsize{10.5bp}{12bp}\selectfont
        #1#2%
        \vspace{2bp}%
    \end{minipage}%
    \hspace{0pt}%
    \vrule width 0.5pt%
    \hspace{4bp}%
    \begin{minipage}[c]{346.32bp}%
        \vspace{2bp}%
        \songti\mdseries\fontsize{10.5bp}{17bp}\selectfont%
        \baselineskip=17bp%
        \setlength{\parindent}{0pt}%
        \color{inkblack}%
        #3%
        \vspace{2bp}%
    \end{minipage}%
    \end{tcolorbox}%
}

% 标签居中版
\newcommand{\pairsection}[2]{%
    \pairsectionbase{\centering}{#1}{#2}%
}

% 标签左对齐版（用于长标签）
\newcommand{\pairsectionleft}[2]{%
    \pairsectionbase{\raggedright\fontsize{10.5bp}{10.5bp}\selectfont}{#1}{#2}%
}


 之前通过
% \newcommand{\fontpath}{...} 覆盖设置
% 默认路径：从 templates/XXX/ 编译时指向 templates/common/fonts/
\providecommand{\fontpath}{../common/fonts/}
\newCJKfontfamily\fzdbs[Path=\fontpath]{FZDBS.ttf}
\newCJKfontfamily\fzssk[Path=\fontpath]{FZSSK.ttf}

% 全局文档颜色设定：公函炭黑
\AtBeginDocument{\color{inkblack}}

% 全局禁止段间距拉伸，确保 tcolorbox 框之间无缝隙
\setlength{\parskip}{0pt}

% ─── 复合标题样式 ───
\newcommand{\lawtitle}[2]{
    \vspace*{0.507cm} 
    \begin{center}
        {\fzdbs\fontsize{22bp}{22bp}\selectfont #1}\par
        \vspace{10.37bp}
        {\fzdbs\fontsize{18bp}{18bp}\selectfont #2}\par
    \end{center}
    \vspace{18.4bp}
}

% ─── 表框（说明框）环境 ───
\newtcolorbox{notesection}[1][]{
    width=467.72bp,
    sharp corners,
    colback=white,
    colframe=inkblack,
    boxrule=0.5pt,
    fontupper=\color{inkblack}\fontsize{10.5bp}{10.5bp}\selectfont,
    enlarge left by=-19.28bp,
    before={\par\nointerlineskip\vspace{-0.5pt}\noindent},
    after={\par},
    boxsep=0pt,
    left=4bp,
    top=4.69bp,
    bottom=4.69bp,
    right=0pt,
    #1
}

% ─── 法律段落换行命令 ───
\newcommand{\lawpar}{\par\vspace*{-4.86bp}}

% ─── 区块副标题样式 ───
\newcommand{\subtitlesection}[1]{%
    \begin{tcolorbox}[
        width=467.72bp,
        sharp corners,
        colback=white,
        colframe=inkblack,
        boxrule=0.5pt,
        boxsep=0pt,
        left=0pt, right=0pt,
        top=4bp, bottom=4bp,
        halign=center,
        fontupper=\fzdbs\fontsize{15bp}{15bp}\selectfont\color{inkblack},
        enlarge left by=-19.28bp,
        before={\par\nointerlineskip\vspace{-0.5pt}\noindent},
        after={\par}
    ]
    #1
    \end{tcolorbox}%
}

% ─── 区块副标题样式（带说明文字） ───
\newcommand{\subtitlenotesection}[2]{%
    \begin{tcolorbox}[
        width=467.72bp, sharp corners, colback=white, colframe=inkblack, boxrule=0.5pt,
        boxsep=0pt, left=0pt, right=0pt, top=4bp, bottom=4bp, halign=center,
        enlarge left by=-19.28bp,
        before={\par\nointerlineskip\vspace{-0.5pt}\noindent}, after={\par}
    ]
    {\fzdbs\fontsize{15bp}{15bp}\selectfont\color{inkblack} #1}\par\nointerlineskip
    \vspace{2bp}
    {\heiti\fontsize{10.5bp}{14bp}\selectfont\color{inkblack} #2}
    \end{tcolorbox}%
}

% ─── 左右分栏基础命令 ───
\newcommand{\pairsectionbase}[3]{%
    \begin{tcolorbox}[
        width=467.72bp, sharp corners, colback=white, colframe=inkblack, boxrule=0.5pt,
        enlarge left by=-19.28bp,
        before={\par\nointerlineskip\vspace{-0.5pt}\noindent}, after={\par},
        boxsep=0pt, left=0pt, right=0pt, top=0pt, bottom=0pt
    ]
    \begin{minipage}[c]{112.9bp}%
        \vspace{2bp}%
        \songti\fontsize{10.5bp}{12bp}\selectfont
        #1#2%
        \vspace{2bp}%
    \end{minipage}%
    \hspace{0pt}%
    \vrule width 0.5pt%
    \hspace{4bp}%
    \begin{minipage}[c]{346.32bp}%
        \vspace{2bp}%
        \songti\mdseries\fontsize{10.5bp}{17bp}\selectfont%
        \baselineskip=17bp%
        \setlength{\parindent}{0pt}%
        \color{inkblack}%
        #3%
        \vspace{2bp}%
    \end{minipage}%
    \end{tcolorbox}%
}

% 标签居中版
\newcommand{\pairsection}[2]{%
    \pairsectionbase{\centering}{#1}{#2}%
}

% 标签左对齐版（用于长标签）
\newcommand{\pairsectionleft}[2]{%
    \pairsectionbase{\raggedright\fontsize{10.5bp}{10.5bp}\selectfont}{#1}{#2}%
}


 之前通过
% \newcommand{\fontpath}{...} 覆盖设置
% 默认路径：从 templates/XXX/ 编译时指向 templates/common/fonts/
\providecommand{\fontpath}{../common/fonts/}
\newCJKfontfamily\fzdbs[Path=\fontpath]{FZDBS.ttf}
\newCJKfontfamily\fzssk[Path=\fontpath]{FZSSK.ttf}

% 全局文档颜色设定：公函炭黑
\AtBeginDocument{\color{inkblack}}

% 全局禁止段间距拉伸，确保 tcolorbox 框之间无缝隙
\setlength{\parskip}{0pt}

% ─── 复合标题样式 ───
\newcommand{\lawtitle}[2]{
    \vspace*{0.507cm} 
    \begin{center}
        {\fzdbs\fontsize{22bp}{22bp}\selectfont #1}\par
        \vspace{10.37bp}
        {\fzdbs\fontsize{18bp}{18bp}\selectfont #2}\par
    \end{center}
    \vspace{18.4bp}
}

% ─── 表框（说明框）环境 ───
\newtcolorbox{notesection}[1][]{
    width=467.72bp,
    sharp corners,
    colback=white,
    colframe=inkblack,
    boxrule=0.5pt,
    fontupper=\color{inkblack}\fontsize{10.5bp}{10.5bp}\selectfont,
    enlarge left by=-19.28bp,
    before={\par\nointerlineskip\vspace{-0.5pt}\noindent},
    after={\par},
    boxsep=0pt,
    left=4bp,
    top=4.69bp,
    bottom=4.69bp,
    right=0pt,
    #1
}

% ─── 法律段落换行命令 ───
\newcommand{\lawpar}{\par\vspace*{-4.86bp}}

% ─── 区块副标题样式 ───
\newcommand{\subtitlesection}[1]{%
    \begin{tcolorbox}[
        width=467.72bp,
        sharp corners,
        colback=white,
        colframe=inkblack,
        boxrule=0.5pt,
        boxsep=0pt,
        left=0pt, right=0pt,
        top=4bp, bottom=4bp,
        halign=center,
        fontupper=\fzdbs\fontsize{15bp}{15bp}\selectfont\color{inkblack},
        enlarge left by=-19.28bp,
        before={\par\nointerlineskip\vspace{-0.5pt}\noindent},
        after={\par}
    ]
    #1
    \end{tcolorbox}%
}

% ─── 区块副标题样式（带说明文字） ───
\newcommand{\subtitlenotesection}[2]{%
    \begin{tcolorbox}[
        width=467.72bp, sharp corners, colback=white, colframe=inkblack, boxrule=0.5pt,
        boxsep=0pt, left=0pt, right=0pt, top=4bp, bottom=4bp, halign=center,
        enlarge left by=-19.28bp,
        before={\par\nointerlineskip\vspace{-0.5pt}\noindent}, after={\par}
    ]
    {\fzdbs\fontsize{15bp}{15bp}\selectfont\color{inkblack} #1}\par\nointerlineskip
    \vspace{2bp}
    {\heiti\fontsize{10.5bp}{14bp}\selectfont\color{inkblack} #2}
    \end{tcolorbox}%
}

% ─── 左右分栏基础命令 ───
\newcommand{\pairsectionbase}[3]{%
    \begin{tcolorbox}[
        width=467.72bp, sharp corners, colback=white, colframe=inkblack, boxrule=0.5pt,
        enlarge left by=-19.28bp,
        before={\par\nointerlineskip\vspace{-0.5pt}\noindent}, after={\par},
        boxsep=0pt, left=0pt, right=0pt, top=0pt, bottom=0pt
    ]
    \begin{minipage}[c]{112.9bp}%
        \vspace{2bp}%
        \songti\fontsize{10.5bp}{12bp}\selectfont
        #1#2%
        \vspace{2bp}%
    \end{minipage}%
    \hspace{0pt}%
    \vrule width 0.5pt%
    \hspace{4bp}%
    \begin{minipage}[c]{346.32bp}%
        \vspace{2bp}%
        \songti\mdseries\fontsize{10.5bp}{17bp}\selectfont%
        \baselineskip=17bp%
        \setlength{\parindent}{0pt}%
        \color{inkblack}%
        #3%
        \vspace{2bp}%
    \end{minipage}%
    \end{tcolorbox}%
}

% 标签居中版
\newcommand{\pairsection}[2]{%
    \pairsectionbase{\centering}{#1}{#2}%
}

% 标签左对齐版（用于长标签）
\newcommand{\pairsectionleft}[2]{%
    \pairsectionbase{\raggedright\fontsize{10.5bp}{10.5bp}\selectfont}{#1}{#2}%
}


 引入
% ═══════════════════════════════════════════════════════════════

\documentclass[a4paper,12pt,twoside]{ctexart}
\usepackage[dvipsnames]{xcolor}
\usepackage{geometry}
\usepackage{tcolorbox}
\usepackage{tikz}
\usepackage{amssymb}
\usepackage{graphicx}
\usepackage{tabularx}
\usepackage{makecell}
\usepackage{multirow}

% 定义公函标准炭黑（K=92%，RGB约0.08）：权威稳重、高质量印刷感
\definecolor{inkblack}{rgb}{0.08,0.08,0.08}

\geometry{a4paper, left=3.18cm, right=3.18cm, top=2.54cm, bottom=2.54cm}

% 配置页脚，实现奇数页右下角，偶数页左下角的页码
\usepackage{fancyhdr}
\pagestyle{fancy}
\fancyhf{}
\renewcommand{\headrulewidth}{0pt}

% 自定义页脚位置微调
\fancyfoot[RO]{\makebox[0pt][l]{\hspace*{28.2pt}\raisebox{2.81pt}{\rmfamily\fontsize{10.5bp}{10.5bp}\selectfont\thepage}}}
\fancyfoot[LE]{\makebox[0pt][r]{\hspace*{33.45pt}\raisebox{2.81pt}{\rmfamily\fontsize{10.5bp}{10.5bp}\selectfont\thepage}}}

% 字体声明
% \fontpath 可由各模板在 % ═══════════════════════════════════════════════════════════════
% 民事起诉状 · 共享前导码
% 所有民事起诉状模板共用本文件，通过 % ═══════════════════════════════════════════════════════════════
% 民事起诉状 · 共享前导码
% 所有民事起诉状模板共用本文件，通过 % ═══════════════════════════════════════════════════════════════
% 民事起诉状 · 共享前导码
% 所有民事起诉状模板共用本文件，通过 \input{../common/preamble} 引入
% ═══════════════════════════════════════════════════════════════

\documentclass[a4paper,12pt,twoside]{ctexart}
\usepackage[dvipsnames]{xcolor}
\usepackage{geometry}
\usepackage{tcolorbox}
\usepackage{tikz}
\usepackage{amssymb}
\usepackage{graphicx}
\usepackage{tabularx}
\usepackage{makecell}
\usepackage{multirow}

% 定义公函标准炭黑（K=92%，RGB约0.08）：权威稳重、高质量印刷感
\definecolor{inkblack}{rgb}{0.08,0.08,0.08}

\geometry{a4paper, left=3.18cm, right=3.18cm, top=2.54cm, bottom=2.54cm}

% 配置页脚，实现奇数页右下角，偶数页左下角的页码
\usepackage{fancyhdr}
\pagestyle{fancy}
\fancyhf{}
\renewcommand{\headrulewidth}{0pt}

% 自定义页脚位置微调
\fancyfoot[RO]{\makebox[0pt][l]{\hspace*{28.2pt}\raisebox{2.81pt}{\rmfamily\fontsize{10.5bp}{10.5bp}\selectfont\thepage}}}
\fancyfoot[LE]{\makebox[0pt][r]{\hspace*{33.45pt}\raisebox{2.81pt}{\rmfamily\fontsize{10.5bp}{10.5bp}\selectfont\thepage}}}

% 字体声明
% \fontpath 可由各模板在 \input{preamble} 之前通过
% \newcommand{\fontpath}{...} 覆盖设置
% 默认路径：从 templates/XXX/ 编译时指向 templates/common/fonts/
\providecommand{\fontpath}{../common/fonts/}
\newCJKfontfamily\fzdbs[Path=\fontpath]{FZDBS.ttf}
\newCJKfontfamily\fzssk[Path=\fontpath]{FZSSK.ttf}

% 全局文档颜色设定：公函炭黑
\AtBeginDocument{\color{inkblack}}

% 全局禁止段间距拉伸，确保 tcolorbox 框之间无缝隙
\setlength{\parskip}{0pt}

% ─── 复合标题样式 ───
\newcommand{\lawtitle}[2]{
    \vspace*{0.507cm} 
    \begin{center}
        {\fzdbs\fontsize{22bp}{22bp}\selectfont #1}\par
        \vspace{10.37bp}
        {\fzdbs\fontsize{18bp}{18bp}\selectfont #2}\par
    \end{center}
    \vspace{18.4bp}
}

% ─── 表框（说明框）环境 ───
\newtcolorbox{notesection}[1][]{
    width=467.72bp,
    sharp corners,
    colback=white,
    colframe=inkblack,
    boxrule=0.5pt,
    fontupper=\color{inkblack}\fontsize{10.5bp}{10.5bp}\selectfont,
    enlarge left by=-19.28bp,
    before={\par\nointerlineskip\vspace{-0.5pt}\noindent},
    after={\par},
    boxsep=0pt,
    left=4bp,
    top=4.69bp,
    bottom=4.69bp,
    right=0pt,
    #1
}

% ─── 法律段落换行命令 ───
\newcommand{\lawpar}{\par\vspace*{-4.86bp}}

% ─── 区块副标题样式 ───
\newcommand{\subtitlesection}[1]{%
    \begin{tcolorbox}[
        width=467.72bp,
        sharp corners,
        colback=white,
        colframe=inkblack,
        boxrule=0.5pt,
        boxsep=0pt,
        left=0pt, right=0pt,
        top=4bp, bottom=4bp,
        halign=center,
        fontupper=\fzdbs\fontsize{15bp}{15bp}\selectfont\color{inkblack},
        enlarge left by=-19.28bp,
        before={\par\nointerlineskip\vspace{-0.5pt}\noindent},
        after={\par}
    ]
    #1
    \end{tcolorbox}%
}

% ─── 区块副标题样式（带说明文字） ───
\newcommand{\subtitlenotesection}[2]{%
    \begin{tcolorbox}[
        width=467.72bp, sharp corners, colback=white, colframe=inkblack, boxrule=0.5pt,
        boxsep=0pt, left=0pt, right=0pt, top=4bp, bottom=4bp, halign=center,
        enlarge left by=-19.28bp,
        before={\par\nointerlineskip\vspace{-0.5pt}\noindent}, after={\par}
    ]
    {\fzdbs\fontsize{15bp}{15bp}\selectfont\color{inkblack} #1}\par\nointerlineskip
    \vspace{2bp}
    {\heiti\fontsize{10.5bp}{14bp}\selectfont\color{inkblack} #2}
    \end{tcolorbox}%
}

% ─── 左右分栏基础命令 ───
\newcommand{\pairsectionbase}[3]{%
    \begin{tcolorbox}[
        width=467.72bp, sharp corners, colback=white, colframe=inkblack, boxrule=0.5pt,
        enlarge left by=-19.28bp,
        before={\par\nointerlineskip\vspace{-0.5pt}\noindent}, after={\par},
        boxsep=0pt, left=0pt, right=0pt, top=0pt, bottom=0pt
    ]
    \begin{minipage}[c]{112.9bp}%
        \vspace{2bp}%
        \songti\fontsize{10.5bp}{12bp}\selectfont
        #1#2%
        \vspace{2bp}%
    \end{minipage}%
    \hspace{0pt}%
    \vrule width 0.5pt%
    \hspace{4bp}%
    \begin{minipage}[c]{346.32bp}%
        \vspace{2bp}%
        \songti\mdseries\fontsize{10.5bp}{17bp}\selectfont%
        \baselineskip=17bp%
        \setlength{\parindent}{0pt}%
        \color{inkblack}%
        #3%
        \vspace{2bp}%
    \end{minipage}%
    \end{tcolorbox}%
}

% 标签居中版
\newcommand{\pairsection}[2]{%
    \pairsectionbase{\centering}{#1}{#2}%
}

% 标签左对齐版（用于长标签）
\newcommand{\pairsectionleft}[2]{%
    \pairsectionbase{\raggedright\fontsize{10.5bp}{10.5bp}\selectfont}{#1}{#2}%
}


 引入
% ═══════════════════════════════════════════════════════════════

\documentclass[a4paper,12pt,twoside]{ctexart}
\usepackage[dvipsnames]{xcolor}
\usepackage{geometry}
\usepackage{tcolorbox}
\usepackage{tikz}
\usepackage{amssymb}
\usepackage{graphicx}
\usepackage{tabularx}
\usepackage{makecell}
\usepackage{multirow}

% 定义公函标准炭黑（K=92%，RGB约0.08）：权威稳重、高质量印刷感
\definecolor{inkblack}{rgb}{0.08,0.08,0.08}

\geometry{a4paper, left=3.18cm, right=3.18cm, top=2.54cm, bottom=2.54cm}

% 配置页脚，实现奇数页右下角，偶数页左下角的页码
\usepackage{fancyhdr}
\pagestyle{fancy}
\fancyhf{}
\renewcommand{\headrulewidth}{0pt}

% 自定义页脚位置微调
\fancyfoot[RO]{\makebox[0pt][l]{\hspace*{28.2pt}\raisebox{2.81pt}{\rmfamily\fontsize{10.5bp}{10.5bp}\selectfont\thepage}}}
\fancyfoot[LE]{\makebox[0pt][r]{\hspace*{33.45pt}\raisebox{2.81pt}{\rmfamily\fontsize{10.5bp}{10.5bp}\selectfont\thepage}}}

% 字体声明
% \fontpath 可由各模板在 % ═══════════════════════════════════════════════════════════════
% 民事起诉状 · 共享前导码
% 所有民事起诉状模板共用本文件，通过 \input{../common/preamble} 引入
% ═══════════════════════════════════════════════════════════════

\documentclass[a4paper,12pt,twoside]{ctexart}
\usepackage[dvipsnames]{xcolor}
\usepackage{geometry}
\usepackage{tcolorbox}
\usepackage{tikz}
\usepackage{amssymb}
\usepackage{graphicx}
\usepackage{tabularx}
\usepackage{makecell}
\usepackage{multirow}

% 定义公函标准炭黑（K=92%，RGB约0.08）：权威稳重、高质量印刷感
\definecolor{inkblack}{rgb}{0.08,0.08,0.08}

\geometry{a4paper, left=3.18cm, right=3.18cm, top=2.54cm, bottom=2.54cm}

% 配置页脚，实现奇数页右下角，偶数页左下角的页码
\usepackage{fancyhdr}
\pagestyle{fancy}
\fancyhf{}
\renewcommand{\headrulewidth}{0pt}

% 自定义页脚位置微调
\fancyfoot[RO]{\makebox[0pt][l]{\hspace*{28.2pt}\raisebox{2.81pt}{\rmfamily\fontsize{10.5bp}{10.5bp}\selectfont\thepage}}}
\fancyfoot[LE]{\makebox[0pt][r]{\hspace*{33.45pt}\raisebox{2.81pt}{\rmfamily\fontsize{10.5bp}{10.5bp}\selectfont\thepage}}}

% 字体声明
% \fontpath 可由各模板在 \input{preamble} 之前通过
% \newcommand{\fontpath}{...} 覆盖设置
% 默认路径：从 templates/XXX/ 编译时指向 templates/common/fonts/
\providecommand{\fontpath}{../common/fonts/}
\newCJKfontfamily\fzdbs[Path=\fontpath]{FZDBS.ttf}
\newCJKfontfamily\fzssk[Path=\fontpath]{FZSSK.ttf}

% 全局文档颜色设定：公函炭黑
\AtBeginDocument{\color{inkblack}}

% 全局禁止段间距拉伸，确保 tcolorbox 框之间无缝隙
\setlength{\parskip}{0pt}

% ─── 复合标题样式 ───
\newcommand{\lawtitle}[2]{
    \vspace*{0.507cm} 
    \begin{center}
        {\fzdbs\fontsize{22bp}{22bp}\selectfont #1}\par
        \vspace{10.37bp}
        {\fzdbs\fontsize{18bp}{18bp}\selectfont #2}\par
    \end{center}
    \vspace{18.4bp}
}

% ─── 表框（说明框）环境 ───
\newtcolorbox{notesection}[1][]{
    width=467.72bp,
    sharp corners,
    colback=white,
    colframe=inkblack,
    boxrule=0.5pt,
    fontupper=\color{inkblack}\fontsize{10.5bp}{10.5bp}\selectfont,
    enlarge left by=-19.28bp,
    before={\par\nointerlineskip\vspace{-0.5pt}\noindent},
    after={\par},
    boxsep=0pt,
    left=4bp,
    top=4.69bp,
    bottom=4.69bp,
    right=0pt,
    #1
}

% ─── 法律段落换行命令 ───
\newcommand{\lawpar}{\par\vspace*{-4.86bp}}

% ─── 区块副标题样式 ───
\newcommand{\subtitlesection}[1]{%
    \begin{tcolorbox}[
        width=467.72bp,
        sharp corners,
        colback=white,
        colframe=inkblack,
        boxrule=0.5pt,
        boxsep=0pt,
        left=0pt, right=0pt,
        top=4bp, bottom=4bp,
        halign=center,
        fontupper=\fzdbs\fontsize{15bp}{15bp}\selectfont\color{inkblack},
        enlarge left by=-19.28bp,
        before={\par\nointerlineskip\vspace{-0.5pt}\noindent},
        after={\par}
    ]
    #1
    \end{tcolorbox}%
}

% ─── 区块副标题样式（带说明文字） ───
\newcommand{\subtitlenotesection}[2]{%
    \begin{tcolorbox}[
        width=467.72bp, sharp corners, colback=white, colframe=inkblack, boxrule=0.5pt,
        boxsep=0pt, left=0pt, right=0pt, top=4bp, bottom=4bp, halign=center,
        enlarge left by=-19.28bp,
        before={\par\nointerlineskip\vspace{-0.5pt}\noindent}, after={\par}
    ]
    {\fzdbs\fontsize{15bp}{15bp}\selectfont\color{inkblack} #1}\par\nointerlineskip
    \vspace{2bp}
    {\heiti\fontsize{10.5bp}{14bp}\selectfont\color{inkblack} #2}
    \end{tcolorbox}%
}

% ─── 左右分栏基础命令 ───
\newcommand{\pairsectionbase}[3]{%
    \begin{tcolorbox}[
        width=467.72bp, sharp corners, colback=white, colframe=inkblack, boxrule=0.5pt,
        enlarge left by=-19.28bp,
        before={\par\nointerlineskip\vspace{-0.5pt}\noindent}, after={\par},
        boxsep=0pt, left=0pt, right=0pt, top=0pt, bottom=0pt
    ]
    \begin{minipage}[c]{112.9bp}%
        \vspace{2bp}%
        \songti\fontsize{10.5bp}{12bp}\selectfont
        #1#2%
        \vspace{2bp}%
    \end{minipage}%
    \hspace{0pt}%
    \vrule width 0.5pt%
    \hspace{4bp}%
    \begin{minipage}[c]{346.32bp}%
        \vspace{2bp}%
        \songti\mdseries\fontsize{10.5bp}{17bp}\selectfont%
        \baselineskip=17bp%
        \setlength{\parindent}{0pt}%
        \color{inkblack}%
        #3%
        \vspace{2bp}%
    \end{minipage}%
    \end{tcolorbox}%
}

% 标签居中版
\newcommand{\pairsection}[2]{%
    \pairsectionbase{\centering}{#1}{#2}%
}

% 标签左对齐版（用于长标签）
\newcommand{\pairsectionleft}[2]{%
    \pairsectionbase{\raggedright\fontsize{10.5bp}{10.5bp}\selectfont}{#1}{#2}%
}


 之前通过
% \newcommand{\fontpath}{...} 覆盖设置
% 默认路径：从 templates/XXX/ 编译时指向 templates/common/fonts/
\providecommand{\fontpath}{../common/fonts/}
\newCJKfontfamily\fzdbs[Path=\fontpath]{FZDBS.ttf}
\newCJKfontfamily\fzssk[Path=\fontpath]{FZSSK.ttf}

% 全局文档颜色设定：公函炭黑
\AtBeginDocument{\color{inkblack}}

% 全局禁止段间距拉伸，确保 tcolorbox 框之间无缝隙
\setlength{\parskip}{0pt}

% ─── 复合标题样式 ───
\newcommand{\lawtitle}[2]{
    \vspace*{0.507cm} 
    \begin{center}
        {\fzdbs\fontsize{22bp}{22bp}\selectfont #1}\par
        \vspace{10.37bp}
        {\fzdbs\fontsize{18bp}{18bp}\selectfont #2}\par
    \end{center}
    \vspace{18.4bp}
}

% ─── 表框（说明框）环境 ───
\newtcolorbox{notesection}[1][]{
    width=467.72bp,
    sharp corners,
    colback=white,
    colframe=inkblack,
    boxrule=0.5pt,
    fontupper=\color{inkblack}\fontsize{10.5bp}{10.5bp}\selectfont,
    enlarge left by=-19.28bp,
    before={\par\nointerlineskip\vspace{-0.5pt}\noindent},
    after={\par},
    boxsep=0pt,
    left=4bp,
    top=4.69bp,
    bottom=4.69bp,
    right=0pt,
    #1
}

% ─── 法律段落换行命令 ───
\newcommand{\lawpar}{\par\vspace*{-4.86bp}}

% ─── 区块副标题样式 ───
\newcommand{\subtitlesection}[1]{%
    \begin{tcolorbox}[
        width=467.72bp,
        sharp corners,
        colback=white,
        colframe=inkblack,
        boxrule=0.5pt,
        boxsep=0pt,
        left=0pt, right=0pt,
        top=4bp, bottom=4bp,
        halign=center,
        fontupper=\fzdbs\fontsize{15bp}{15bp}\selectfont\color{inkblack},
        enlarge left by=-19.28bp,
        before={\par\nointerlineskip\vspace{-0.5pt}\noindent},
        after={\par}
    ]
    #1
    \end{tcolorbox}%
}

% ─── 区块副标题样式（带说明文字） ───
\newcommand{\subtitlenotesection}[2]{%
    \begin{tcolorbox}[
        width=467.72bp, sharp corners, colback=white, colframe=inkblack, boxrule=0.5pt,
        boxsep=0pt, left=0pt, right=0pt, top=4bp, bottom=4bp, halign=center,
        enlarge left by=-19.28bp,
        before={\par\nointerlineskip\vspace{-0.5pt}\noindent}, after={\par}
    ]
    {\fzdbs\fontsize{15bp}{15bp}\selectfont\color{inkblack} #1}\par\nointerlineskip
    \vspace{2bp}
    {\heiti\fontsize{10.5bp}{14bp}\selectfont\color{inkblack} #2}
    \end{tcolorbox}%
}

% ─── 左右分栏基础命令 ───
\newcommand{\pairsectionbase}[3]{%
    \begin{tcolorbox}[
        width=467.72bp, sharp corners, colback=white, colframe=inkblack, boxrule=0.5pt,
        enlarge left by=-19.28bp,
        before={\par\nointerlineskip\vspace{-0.5pt}\noindent}, after={\par},
        boxsep=0pt, left=0pt, right=0pt, top=0pt, bottom=0pt
    ]
    \begin{minipage}[c]{112.9bp}%
        \vspace{2bp}%
        \songti\fontsize{10.5bp}{12bp}\selectfont
        #1#2%
        \vspace{2bp}%
    \end{minipage}%
    \hspace{0pt}%
    \vrule width 0.5pt%
    \hspace{4bp}%
    \begin{minipage}[c]{346.32bp}%
        \vspace{2bp}%
        \songti\mdseries\fontsize{10.5bp}{17bp}\selectfont%
        \baselineskip=17bp%
        \setlength{\parindent}{0pt}%
        \color{inkblack}%
        #3%
        \vspace{2bp}%
    \end{minipage}%
    \end{tcolorbox}%
}

% 标签居中版
\newcommand{\pairsection}[2]{%
    \pairsectionbase{\centering}{#1}{#2}%
}

% 标签左对齐版（用于长标签）
\newcommand{\pairsectionleft}[2]{%
    \pairsectionbase{\raggedright\fontsize{10.5bp}{10.5bp}\selectfont}{#1}{#2}%
}


 引入
% ═══════════════════════════════════════════════════════════════

\documentclass[a4paper,12pt,twoside]{ctexart}
\usepackage[dvipsnames]{xcolor}
\usepackage{geometry}
\usepackage{tcolorbox}
\usepackage{tikz}
\usepackage{amssymb}
\usepackage{graphicx}
\usepackage{tabularx}
\usepackage{makecell}
\usepackage{multirow}

% 定义公函标准炭黑（K=92%，RGB约0.08）：权威稳重、高质量印刷感
\definecolor{inkblack}{rgb}{0.08,0.08,0.08}

\geometry{a4paper, left=3.18cm, right=3.18cm, top=2.54cm, bottom=2.54cm}

% 配置页脚，实现奇数页右下角，偶数页左下角的页码
\usepackage{fancyhdr}
\pagestyle{fancy}
\fancyhf{}
\renewcommand{\headrulewidth}{0pt}

% 自定义页脚位置微调
\fancyfoot[RO]{\makebox[0pt][l]{\hspace*{28.2pt}\raisebox{2.81pt}{\rmfamily\fontsize{10.5bp}{10.5bp}\selectfont\thepage}}}
\fancyfoot[LE]{\makebox[0pt][r]{\hspace*{33.45pt}\raisebox{2.81pt}{\rmfamily\fontsize{10.5bp}{10.5bp}\selectfont\thepage}}}

% 字体声明
% \fontpath 可由各模板在 % ═══════════════════════════════════════════════════════════════
% 民事起诉状 · 共享前导码
% 所有民事起诉状模板共用本文件，通过 % ═══════════════════════════════════════════════════════════════
% 民事起诉状 · 共享前导码
% 所有民事起诉状模板共用本文件，通过 \input{../common/preamble} 引入
% ═══════════════════════════════════════════════════════════════

\documentclass[a4paper,12pt,twoside]{ctexart}
\usepackage[dvipsnames]{xcolor}
\usepackage{geometry}
\usepackage{tcolorbox}
\usepackage{tikz}
\usepackage{amssymb}
\usepackage{graphicx}
\usepackage{tabularx}
\usepackage{makecell}
\usepackage{multirow}

% 定义公函标准炭黑（K=92%，RGB约0.08）：权威稳重、高质量印刷感
\definecolor{inkblack}{rgb}{0.08,0.08,0.08}

\geometry{a4paper, left=3.18cm, right=3.18cm, top=2.54cm, bottom=2.54cm}

% 配置页脚，实现奇数页右下角，偶数页左下角的页码
\usepackage{fancyhdr}
\pagestyle{fancy}
\fancyhf{}
\renewcommand{\headrulewidth}{0pt}

% 自定义页脚位置微调
\fancyfoot[RO]{\makebox[0pt][l]{\hspace*{28.2pt}\raisebox{2.81pt}{\rmfamily\fontsize{10.5bp}{10.5bp}\selectfont\thepage}}}
\fancyfoot[LE]{\makebox[0pt][r]{\hspace*{33.45pt}\raisebox{2.81pt}{\rmfamily\fontsize{10.5bp}{10.5bp}\selectfont\thepage}}}

% 字体声明
% \fontpath 可由各模板在 \input{preamble} 之前通过
% \newcommand{\fontpath}{...} 覆盖设置
% 默认路径：从 templates/XXX/ 编译时指向 templates/common/fonts/
\providecommand{\fontpath}{../common/fonts/}
\newCJKfontfamily\fzdbs[Path=\fontpath]{FZDBS.ttf}
\newCJKfontfamily\fzssk[Path=\fontpath]{FZSSK.ttf}

% 全局文档颜色设定：公函炭黑
\AtBeginDocument{\color{inkblack}}

% 全局禁止段间距拉伸，确保 tcolorbox 框之间无缝隙
\setlength{\parskip}{0pt}

% ─── 复合标题样式 ───
\newcommand{\lawtitle}[2]{
    \vspace*{0.507cm} 
    \begin{center}
        {\fzdbs\fontsize{22bp}{22bp}\selectfont #1}\par
        \vspace{10.37bp}
        {\fzdbs\fontsize{18bp}{18bp}\selectfont #2}\par
    \end{center}
    \vspace{18.4bp}
}

% ─── 表框（说明框）环境 ───
\newtcolorbox{notesection}[1][]{
    width=467.72bp,
    sharp corners,
    colback=white,
    colframe=inkblack,
    boxrule=0.5pt,
    fontupper=\color{inkblack}\fontsize{10.5bp}{10.5bp}\selectfont,
    enlarge left by=-19.28bp,
    before={\par\nointerlineskip\vspace{-0.5pt}\noindent},
    after={\par},
    boxsep=0pt,
    left=4bp,
    top=4.69bp,
    bottom=4.69bp,
    right=0pt,
    #1
}

% ─── 法律段落换行命令 ───
\newcommand{\lawpar}{\par\vspace*{-4.86bp}}

% ─── 区块副标题样式 ───
\newcommand{\subtitlesection}[1]{%
    \begin{tcolorbox}[
        width=467.72bp,
        sharp corners,
        colback=white,
        colframe=inkblack,
        boxrule=0.5pt,
        boxsep=0pt,
        left=0pt, right=0pt,
        top=4bp, bottom=4bp,
        halign=center,
        fontupper=\fzdbs\fontsize{15bp}{15bp}\selectfont\color{inkblack},
        enlarge left by=-19.28bp,
        before={\par\nointerlineskip\vspace{-0.5pt}\noindent},
        after={\par}
    ]
    #1
    \end{tcolorbox}%
}

% ─── 区块副标题样式（带说明文字） ───
\newcommand{\subtitlenotesection}[2]{%
    \begin{tcolorbox}[
        width=467.72bp, sharp corners, colback=white, colframe=inkblack, boxrule=0.5pt,
        boxsep=0pt, left=0pt, right=0pt, top=4bp, bottom=4bp, halign=center,
        enlarge left by=-19.28bp,
        before={\par\nointerlineskip\vspace{-0.5pt}\noindent}, after={\par}
    ]
    {\fzdbs\fontsize{15bp}{15bp}\selectfont\color{inkblack} #1}\par\nointerlineskip
    \vspace{2bp}
    {\heiti\fontsize{10.5bp}{14bp}\selectfont\color{inkblack} #2}
    \end{tcolorbox}%
}

% ─── 左右分栏基础命令 ───
\newcommand{\pairsectionbase}[3]{%
    \begin{tcolorbox}[
        width=467.72bp, sharp corners, colback=white, colframe=inkblack, boxrule=0.5pt,
        enlarge left by=-19.28bp,
        before={\par\nointerlineskip\vspace{-0.5pt}\noindent}, after={\par},
        boxsep=0pt, left=0pt, right=0pt, top=0pt, bottom=0pt
    ]
    \begin{minipage}[c]{112.9bp}%
        \vspace{2bp}%
        \songti\fontsize{10.5bp}{12bp}\selectfont
        #1#2%
        \vspace{2bp}%
    \end{minipage}%
    \hspace{0pt}%
    \vrule width 0.5pt%
    \hspace{4bp}%
    \begin{minipage}[c]{346.32bp}%
        \vspace{2bp}%
        \songti\mdseries\fontsize{10.5bp}{17bp}\selectfont%
        \baselineskip=17bp%
        \setlength{\parindent}{0pt}%
        \color{inkblack}%
        #3%
        \vspace{2bp}%
    \end{minipage}%
    \end{tcolorbox}%
}

% 标签居中版
\newcommand{\pairsection}[2]{%
    \pairsectionbase{\centering}{#1}{#2}%
}

% 标签左对齐版（用于长标签）
\newcommand{\pairsectionleft}[2]{%
    \pairsectionbase{\raggedright\fontsize{10.5bp}{10.5bp}\selectfont}{#1}{#2}%
}


 引入
% ═══════════════════════════════════════════════════════════════

\documentclass[a4paper,12pt,twoside]{ctexart}
\usepackage[dvipsnames]{xcolor}
\usepackage{geometry}
\usepackage{tcolorbox}
\usepackage{tikz}
\usepackage{amssymb}
\usepackage{graphicx}
\usepackage{tabularx}
\usepackage{makecell}
\usepackage{multirow}

% 定义公函标准炭黑（K=92%，RGB约0.08）：权威稳重、高质量印刷感
\definecolor{inkblack}{rgb}{0.08,0.08,0.08}

\geometry{a4paper, left=3.18cm, right=3.18cm, top=2.54cm, bottom=2.54cm}

% 配置页脚，实现奇数页右下角，偶数页左下角的页码
\usepackage{fancyhdr}
\pagestyle{fancy}
\fancyhf{}
\renewcommand{\headrulewidth}{0pt}

% 自定义页脚位置微调
\fancyfoot[RO]{\makebox[0pt][l]{\hspace*{28.2pt}\raisebox{2.81pt}{\rmfamily\fontsize{10.5bp}{10.5bp}\selectfont\thepage}}}
\fancyfoot[LE]{\makebox[0pt][r]{\hspace*{33.45pt}\raisebox{2.81pt}{\rmfamily\fontsize{10.5bp}{10.5bp}\selectfont\thepage}}}

% 字体声明
% \fontpath 可由各模板在 % ═══════════════════════════════════════════════════════════════
% 民事起诉状 · 共享前导码
% 所有民事起诉状模板共用本文件，通过 \input{../common/preamble} 引入
% ═══════════════════════════════════════════════════════════════

\documentclass[a4paper,12pt,twoside]{ctexart}
\usepackage[dvipsnames]{xcolor}
\usepackage{geometry}
\usepackage{tcolorbox}
\usepackage{tikz}
\usepackage{amssymb}
\usepackage{graphicx}
\usepackage{tabularx}
\usepackage{makecell}
\usepackage{multirow}

% 定义公函标准炭黑（K=92%，RGB约0.08）：权威稳重、高质量印刷感
\definecolor{inkblack}{rgb}{0.08,0.08,0.08}

\geometry{a4paper, left=3.18cm, right=3.18cm, top=2.54cm, bottom=2.54cm}

% 配置页脚，实现奇数页右下角，偶数页左下角的页码
\usepackage{fancyhdr}
\pagestyle{fancy}
\fancyhf{}
\renewcommand{\headrulewidth}{0pt}

% 自定义页脚位置微调
\fancyfoot[RO]{\makebox[0pt][l]{\hspace*{28.2pt}\raisebox{2.81pt}{\rmfamily\fontsize{10.5bp}{10.5bp}\selectfont\thepage}}}
\fancyfoot[LE]{\makebox[0pt][r]{\hspace*{33.45pt}\raisebox{2.81pt}{\rmfamily\fontsize{10.5bp}{10.5bp}\selectfont\thepage}}}

% 字体声明
% \fontpath 可由各模板在 \input{preamble} 之前通过
% \newcommand{\fontpath}{...} 覆盖设置
% 默认路径：从 templates/XXX/ 编译时指向 templates/common/fonts/
\providecommand{\fontpath}{../common/fonts/}
\newCJKfontfamily\fzdbs[Path=\fontpath]{FZDBS.ttf}
\newCJKfontfamily\fzssk[Path=\fontpath]{FZSSK.ttf}

% 全局文档颜色设定：公函炭黑
\AtBeginDocument{\color{inkblack}}

% 全局禁止段间距拉伸，确保 tcolorbox 框之间无缝隙
\setlength{\parskip}{0pt}

% ─── 复合标题样式 ───
\newcommand{\lawtitle}[2]{
    \vspace*{0.507cm} 
    \begin{center}
        {\fzdbs\fontsize{22bp}{22bp}\selectfont #1}\par
        \vspace{10.37bp}
        {\fzdbs\fontsize{18bp}{18bp}\selectfont #2}\par
    \end{center}
    \vspace{18.4bp}
}

% ─── 表框（说明框）环境 ───
\newtcolorbox{notesection}[1][]{
    width=467.72bp,
    sharp corners,
    colback=white,
    colframe=inkblack,
    boxrule=0.5pt,
    fontupper=\color{inkblack}\fontsize{10.5bp}{10.5bp}\selectfont,
    enlarge left by=-19.28bp,
    before={\par\nointerlineskip\vspace{-0.5pt}\noindent},
    after={\par},
    boxsep=0pt,
    left=4bp,
    top=4.69bp,
    bottom=4.69bp,
    right=0pt,
    #1
}

% ─── 法律段落换行命令 ───
\newcommand{\lawpar}{\par\vspace*{-4.86bp}}

% ─── 区块副标题样式 ───
\newcommand{\subtitlesection}[1]{%
    \begin{tcolorbox}[
        width=467.72bp,
        sharp corners,
        colback=white,
        colframe=inkblack,
        boxrule=0.5pt,
        boxsep=0pt,
        left=0pt, right=0pt,
        top=4bp, bottom=4bp,
        halign=center,
        fontupper=\fzdbs\fontsize{15bp}{15bp}\selectfont\color{inkblack},
        enlarge left by=-19.28bp,
        before={\par\nointerlineskip\vspace{-0.5pt}\noindent},
        after={\par}
    ]
    #1
    \end{tcolorbox}%
}

% ─── 区块副标题样式（带说明文字） ───
\newcommand{\subtitlenotesection}[2]{%
    \begin{tcolorbox}[
        width=467.72bp, sharp corners, colback=white, colframe=inkblack, boxrule=0.5pt,
        boxsep=0pt, left=0pt, right=0pt, top=4bp, bottom=4bp, halign=center,
        enlarge left by=-19.28bp,
        before={\par\nointerlineskip\vspace{-0.5pt}\noindent}, after={\par}
    ]
    {\fzdbs\fontsize{15bp}{15bp}\selectfont\color{inkblack} #1}\par\nointerlineskip
    \vspace{2bp}
    {\heiti\fontsize{10.5bp}{14bp}\selectfont\color{inkblack} #2}
    \end{tcolorbox}%
}

% ─── 左右分栏基础命令 ───
\newcommand{\pairsectionbase}[3]{%
    \begin{tcolorbox}[
        width=467.72bp, sharp corners, colback=white, colframe=inkblack, boxrule=0.5pt,
        enlarge left by=-19.28bp,
        before={\par\nointerlineskip\vspace{-0.5pt}\noindent}, after={\par},
        boxsep=0pt, left=0pt, right=0pt, top=0pt, bottom=0pt
    ]
    \begin{minipage}[c]{112.9bp}%
        \vspace{2bp}%
        \songti\fontsize{10.5bp}{12bp}\selectfont
        #1#2%
        \vspace{2bp}%
    \end{minipage}%
    \hspace{0pt}%
    \vrule width 0.5pt%
    \hspace{4bp}%
    \begin{minipage}[c]{346.32bp}%
        \vspace{2bp}%
        \songti\mdseries\fontsize{10.5bp}{17bp}\selectfont%
        \baselineskip=17bp%
        \setlength{\parindent}{0pt}%
        \color{inkblack}%
        #3%
        \vspace{2bp}%
    \end{minipage}%
    \end{tcolorbox}%
}

% 标签居中版
\newcommand{\pairsection}[2]{%
    \pairsectionbase{\centering}{#1}{#2}%
}

% 标签左对齐版（用于长标签）
\newcommand{\pairsectionleft}[2]{%
    \pairsectionbase{\raggedright\fontsize{10.5bp}{10.5bp}\selectfont}{#1}{#2}%
}


 之前通过
% \newcommand{\fontpath}{...} 覆盖设置
% 默认路径：从 templates/XXX/ 编译时指向 templates/common/fonts/
\providecommand{\fontpath}{../common/fonts/}
\newCJKfontfamily\fzdbs[Path=\fontpath]{FZDBS.ttf}
\newCJKfontfamily\fzssk[Path=\fontpath]{FZSSK.ttf}

% 全局文档颜色设定：公函炭黑
\AtBeginDocument{\color{inkblack}}

% 全局禁止段间距拉伸，确保 tcolorbox 框之间无缝隙
\setlength{\parskip}{0pt}

% ─── 复合标题样式 ───
\newcommand{\lawtitle}[2]{
    \vspace*{0.507cm} 
    \begin{center}
        {\fzdbs\fontsize{22bp}{22bp}\selectfont #1}\par
        \vspace{10.37bp}
        {\fzdbs\fontsize{18bp}{18bp}\selectfont #2}\par
    \end{center}
    \vspace{18.4bp}
}

% ─── 表框（说明框）环境 ───
\newtcolorbox{notesection}[1][]{
    width=467.72bp,
    sharp corners,
    colback=white,
    colframe=inkblack,
    boxrule=0.5pt,
    fontupper=\color{inkblack}\fontsize{10.5bp}{10.5bp}\selectfont,
    enlarge left by=-19.28bp,
    before={\par\nointerlineskip\vspace{-0.5pt}\noindent},
    after={\par},
    boxsep=0pt,
    left=4bp,
    top=4.69bp,
    bottom=4.69bp,
    right=0pt,
    #1
}

% ─── 法律段落换行命令 ───
\newcommand{\lawpar}{\par\vspace*{-4.86bp}}

% ─── 区块副标题样式 ───
\newcommand{\subtitlesection}[1]{%
    \begin{tcolorbox}[
        width=467.72bp,
        sharp corners,
        colback=white,
        colframe=inkblack,
        boxrule=0.5pt,
        boxsep=0pt,
        left=0pt, right=0pt,
        top=4bp, bottom=4bp,
        halign=center,
        fontupper=\fzdbs\fontsize{15bp}{15bp}\selectfont\color{inkblack},
        enlarge left by=-19.28bp,
        before={\par\nointerlineskip\vspace{-0.5pt}\noindent},
        after={\par}
    ]
    #1
    \end{tcolorbox}%
}

% ─── 区块副标题样式（带说明文字） ───
\newcommand{\subtitlenotesection}[2]{%
    \begin{tcolorbox}[
        width=467.72bp, sharp corners, colback=white, colframe=inkblack, boxrule=0.5pt,
        boxsep=0pt, left=0pt, right=0pt, top=4bp, bottom=4bp, halign=center,
        enlarge left by=-19.28bp,
        before={\par\nointerlineskip\vspace{-0.5pt}\noindent}, after={\par}
    ]
    {\fzdbs\fontsize{15bp}{15bp}\selectfont\color{inkblack} #1}\par\nointerlineskip
    \vspace{2bp}
    {\heiti\fontsize{10.5bp}{14bp}\selectfont\color{inkblack} #2}
    \end{tcolorbox}%
}

% ─── 左右分栏基础命令 ───
\newcommand{\pairsectionbase}[3]{%
    \begin{tcolorbox}[
        width=467.72bp, sharp corners, colback=white, colframe=inkblack, boxrule=0.5pt,
        enlarge left by=-19.28bp,
        before={\par\nointerlineskip\vspace{-0.5pt}\noindent}, after={\par},
        boxsep=0pt, left=0pt, right=0pt, top=0pt, bottom=0pt
    ]
    \begin{minipage}[c]{112.9bp}%
        \vspace{2bp}%
        \songti\fontsize{10.5bp}{12bp}\selectfont
        #1#2%
        \vspace{2bp}%
    \end{minipage}%
    \hspace{0pt}%
    \vrule width 0.5pt%
    \hspace{4bp}%
    \begin{minipage}[c]{346.32bp}%
        \vspace{2bp}%
        \songti\mdseries\fontsize{10.5bp}{17bp}\selectfont%
        \baselineskip=17bp%
        \setlength{\parindent}{0pt}%
        \color{inkblack}%
        #3%
        \vspace{2bp}%
    \end{minipage}%
    \end{tcolorbox}%
}

% 标签居中版
\newcommand{\pairsection}[2]{%
    \pairsectionbase{\centering}{#1}{#2}%
}

% 标签左对齐版（用于长标签）
\newcommand{\pairsectionleft}[2]{%
    \pairsectionbase{\raggedright\fontsize{10.5bp}{10.5bp}\selectfont}{#1}{#2}%
}


 之前通过
% \newcommand{\fontpath}{...} 覆盖设置
% 默认路径：从 templates/XXX/ 编译时指向 templates/common/fonts/
\providecommand{\fontpath}{../common/fonts/}
\newCJKfontfamily\fzdbs[Path=\fontpath]{FZDBS.ttf}
\newCJKfontfamily\fzssk[Path=\fontpath]{FZSSK.ttf}

% 全局文档颜色设定：公函炭黑
\AtBeginDocument{\color{inkblack}}

% 全局禁止段间距拉伸，确保 tcolorbox 框之间无缝隙
\setlength{\parskip}{0pt}

% ─── 复合标题样式 ───
\newcommand{\lawtitle}[2]{
    \vspace*{0.507cm} 
    \begin{center}
        {\fzdbs\fontsize{22bp}{22bp}\selectfont #1}\par
        \vspace{10.37bp}
        {\fzdbs\fontsize{18bp}{18bp}\selectfont #2}\par
    \end{center}
    \vspace{18.4bp}
}

% ─── 表框（说明框）环境 ───
\newtcolorbox{notesection}[1][]{
    width=467.72bp,
    sharp corners,
    colback=white,
    colframe=inkblack,
    boxrule=0.5pt,
    fontupper=\color{inkblack}\fontsize{10.5bp}{10.5bp}\selectfont,
    enlarge left by=-19.28bp,
    before={\par\nointerlineskip\vspace{-0.5pt}\noindent},
    after={\par},
    boxsep=0pt,
    left=4bp,
    top=4.69bp,
    bottom=4.69bp,
    right=0pt,
    #1
}

% ─── 法律段落换行命令 ───
\newcommand{\lawpar}{\par\vspace*{-4.86bp}}

% ─── 区块副标题样式 ───
\newcommand{\subtitlesection}[1]{%
    \begin{tcolorbox}[
        width=467.72bp,
        sharp corners,
        colback=white,
        colframe=inkblack,
        boxrule=0.5pt,
        boxsep=0pt,
        left=0pt, right=0pt,
        top=4bp, bottom=4bp,
        halign=center,
        fontupper=\fzdbs\fontsize{15bp}{15bp}\selectfont\color{inkblack},
        enlarge left by=-19.28bp,
        before={\par\nointerlineskip\vspace{-0.5pt}\noindent},
        after={\par}
    ]
    #1
    \end{tcolorbox}%
}

% ─── 区块副标题样式（带说明文字） ───
\newcommand{\subtitlenotesection}[2]{%
    \begin{tcolorbox}[
        width=467.72bp, sharp corners, colback=white, colframe=inkblack, boxrule=0.5pt,
        boxsep=0pt, left=0pt, right=0pt, top=4bp, bottom=4bp, halign=center,
        enlarge left by=-19.28bp,
        before={\par\nointerlineskip\vspace{-0.5pt}\noindent}, after={\par}
    ]
    {\fzdbs\fontsize{15bp}{15bp}\selectfont\color{inkblack} #1}\par\nointerlineskip
    \vspace{2bp}
    {\heiti\fontsize{10.5bp}{14bp}\selectfont\color{inkblack} #2}
    \end{tcolorbox}%
}

% ─── 左右分栏基础命令 ───
\newcommand{\pairsectionbase}[3]{%
    \begin{tcolorbox}[
        width=467.72bp, sharp corners, colback=white, colframe=inkblack, boxrule=0.5pt,
        enlarge left by=-19.28bp,
        before={\par\nointerlineskip\vspace{-0.5pt}\noindent}, after={\par},
        boxsep=0pt, left=0pt, right=0pt, top=0pt, bottom=0pt
    ]
    \begin{minipage}[c]{112.9bp}%
        \vspace{2bp}%
        \songti\fontsize{10.5bp}{12bp}\selectfont
        #1#2%
        \vspace{2bp}%
    \end{minipage}%
    \hspace{0pt}%
    \vrule width 0.5pt%
    \hspace{4bp}%
    \begin{minipage}[c]{346.32bp}%
        \vspace{2bp}%
        \songti\mdseries\fontsize{10.5bp}{17bp}\selectfont%
        \baselineskip=17bp%
        \setlength{\parindent}{0pt}%
        \color{inkblack}%
        #3%
        \vspace{2bp}%
    \end{minipage}%
    \end{tcolorbox}%
}

% 标签居中版
\newcommand{\pairsection}[2]{%
    \pairsectionbase{\centering}{#1}{#2}%
}

% 标签左对齐版（用于长标签）
\newcommand{\pairsectionleft}[2]{%
    \pairsectionbase{\raggedright\fontsize{10.5bp}{10.5bp}\selectfont}{#1}{#2}%
}


 之前通过
% \newcommand{\fontpath}{...} 覆盖设置
% 默认路径：从 templates/XXX/ 编译时指向 templates/common/fonts/
\providecommand{\fontpath}{../common/fonts/}
\newCJKfontfamily\fzdbs[Path=\fontpath]{FZDBS.ttf}
\newCJKfontfamily\fzssk[Path=\fontpath]{FZSSK.ttf}

% 全局文档颜色设定：公函炭黑
\AtBeginDocument{\color{inkblack}}

% 全局禁止段间距拉伸，确保 tcolorbox 框之间无缝隙
\setlength{\parskip}{0pt}

% ─── 复合标题样式 ───
\newcommand{\lawtitle}[2]{
    \vspace*{0.507cm} 
    \begin{center}
        {\fzdbs\fontsize{22bp}{22bp}\selectfont #1}\par
        \vspace{10.37bp}
        {\fzdbs\fontsize{18bp}{18bp}\selectfont #2}\par
    \end{center}
    \vspace{18.4bp}
}

% ─── 表框（说明框）环境 ───
\newtcolorbox{notesection}[1][]{
    width=467.72bp,
    sharp corners,
    colback=white,
    colframe=inkblack,
    boxrule=0.5pt,
    fontupper=\color{inkblack}\fontsize{10.5bp}{10.5bp}\selectfont,
    enlarge left by=-19.28bp,
    before={\par\nointerlineskip\vspace{-0.5pt}\noindent},
    after={\par},
    boxsep=0pt,
    left=4bp,
    top=4.69bp,
    bottom=4.69bp,
    right=0pt,
    #1
}

% ─── 法律段落换行命令 ───
\newcommand{\lawpar}{\par\vspace*{-4.86bp}}

% ─── 区块副标题样式 ───
\newcommand{\subtitlesection}[1]{%
    \begin{tcolorbox}[
        width=467.72bp,
        sharp corners,
        colback=white,
        colframe=inkblack,
        boxrule=0.5pt,
        boxsep=0pt,
        left=0pt, right=0pt,
        top=4bp, bottom=4bp,
        halign=center,
        fontupper=\fzdbs\fontsize{15bp}{15bp}\selectfont\color{inkblack},
        enlarge left by=-19.28bp,
        before={\par\nointerlineskip\vspace{-0.5pt}\noindent},
        after={\par}
    ]
    #1
    \end{tcolorbox}%
}

% ─── 区块副标题样式（带说明文字） ───
\newcommand{\subtitlenotesection}[2]{%
    \begin{tcolorbox}[
        width=467.72bp, sharp corners, colback=white, colframe=inkblack, boxrule=0.5pt,
        boxsep=0pt, left=0pt, right=0pt, top=4bp, bottom=4bp, halign=center,
        enlarge left by=-19.28bp,
        before={\par\nointerlineskip\vspace{-0.5pt}\noindent}, after={\par}
    ]
    {\fzdbs\fontsize{15bp}{15bp}\selectfont\color{inkblack} #1}\par\nointerlineskip
    \vspace{2bp}
    {\heiti\fontsize{10.5bp}{14bp}\selectfont\color{inkblack} #2}
    \end{tcolorbox}%
}

% ─── 左右分栏基础命令 ───
\newcommand{\pairsectionbase}[3]{%
    \begin{tcolorbox}[
        width=467.72bp, sharp corners, colback=white, colframe=inkblack, boxrule=0.5pt,
        enlarge left by=-19.28bp,
        before={\par\nointerlineskip\vspace{-0.5pt}\noindent}, after={\par},
        boxsep=0pt, left=0pt, right=0pt, top=0pt, bottom=0pt
    ]
    \begin{minipage}[c]{112.9bp}%
        \vspace{2bp}%
        \songti\fontsize{10.5bp}{12bp}\selectfont
        #1#2%
        \vspace{2bp}%
    \end{minipage}%
    \hspace{0pt}%
    \vrule width 0.5pt%
    \hspace{4bp}%
    \begin{minipage}[c]{346.32bp}%
        \vspace{2bp}%
        \songti\mdseries\fontsize{10.5bp}{17bp}\selectfont%
        \baselineskip=17bp%
        \setlength{\parindent}{0pt}%
        \color{inkblack}%
        #3%
        \vspace{2bp}%
    \end{minipage}%
    \end{tcolorbox}%
}

% 标签居中版
\newcommand{\pairsection}[2]{%
    \pairsectionbase{\centering}{#1}{#2}%
}

% 标签左对齐版（用于长标签）
\newcommand{\pairsectionleft}[2]{%
    \pairsectionbase{\raggedright\fontsize{10.5bp}{10.5bp}\selectfont}{#1}{#2}%
}


